% -*- coding: utf-8 -*-

\chapter{内容要求（引言或绪论）} \label{chpt:intro}

毕业论文（设计）一般应有前言和文献综述，前言说明论文（设计）的工作目的（背景），并引出课题；文献综述要对国内外研究现状和相关领域中已有研究成果进行简要评述，并介绍毕业论文（设计）中的研究工作设想、方法等内容。

\chapter{格式要求} \label{chpt:format}

各部分要求如下：

\section{封面}

使用教务处统一制作的封面，见本文档给出的封面。

\section{声明}

“关于南开大学本科生毕业论文（设计）的声明”用三号字、黑体，居中书写，正文字体为小四号宋体。学生和指导教师应手写签字，不得打印，不得用印章。声明时间与封面中论文完成时间一致，应在答辩之前。

\section{中文摘要}

关键词是为了便于做文献索引和检索工作而从论文中选取出来用以表示全文主题内容信息的单词或术语，\textcolor{red}{摘要内容后另起一行标明，一般 3～5 个，之间用分号分开。注意：小四号字、宋体。}

\section{英文摘要}

\textcolor{red}{英文摘要另起一页，与中文摘要相对，小四号、Times New Roman。}

\section{目录}

目录应将文内的章节标题依次排列。

\section{正文}

正文是学位论文的主体和核心部分，一般包括以下几个方面：

（1）学位论文字数：按照学校要求，理科类毕业论文（设计）应根据学科特点做出具体要求，其学生毕业论文（设计）正文字数一般不应少于 8000 字（包括图表在内，附录除外）。对于个别基础学科纯理论研究型专业，毕业论文字数一般不应少于 6000 字。\textcolor{red}{软件工程专业毕业论文（设计）字数要求不少于 8000 字。如果是工程属性题目类似“系统设计与实现”的论文，正文一般不少于 40 页；如果是工学属性题目属于理论类，正文不能少于 30 页。此处正文是指从“第一章 绪论”开始到“最后一章 总结和展望”结束之间的部分，不包括论文首页、中英文摘要、目录、参考文献以及致谢部分。}

（2）引言（第一章）：包括研究的目的和意义，问题的提出，选题的背景，文献综述，研究方法，论文结构安排等。

（3）各具体章节：本部分是论文作者的研究内容，是论文的核心。各章之间互相关联，符合逻辑顺序。

（4）总结与展望（最后一章）：是学位论文最终和总体的结论，应明确、精练、完整、准确。着重阐述作者的创造性工作及所取得的研究成果在本学术领域的地位、作用和意义，还可进一步提出需要讨论的问题和建议。

\section{参考文献}

为了反映论文的科学依据和作者尊重他人研究成果的严肃态度以及向读者提供有关信息的出处，应列出参考文献表。参考文献表中列出的一般应限于作者直接阅读过的、最主要的、发表在正式出版物上的文献。私人通信和未公开发表的资料，一般不宜列入参考文献，可紧跟在引用的内容之后注释或标注在当页的下方。\textcolor{red}{参考文献视论文写作实际需要而定，原则上要求参考文献不少于 15 篇，其中外文文献 2 篇以上，全部在正文中被引用。鼓励学生引用近五年的最新学术研究成果。}

《中国高校自然科学学报编排规范》中几种主要参考文献著录表的格式为：

连续出版物：作者.文题.刊名，年，卷号（期号）：起~止页码。

专（译）著：作者.书名（译者）.出版地：出版者，出版年，起~止页码。

论文集：作者.文题.见（In）：编者，编（eds.）文集名.出版地：出版者，出版年，起~止页码。

学位论文：作者.文题:[博士或硕士学位论文].授予单位，授予年。

专\hspace{2em}利：申请者.专利名.国家.专利文献种类.专利号，授权日期。

技术标准：发布单位.技术标准代号.技术标准名称.出版地：出版者，出版日期。

人文社会科学论文的文献资料格式为：

图书：著者.书名.出版者，出版时间，版次，页次。

期刊：作者.篇名.期刊名称，期号。

报纸：作者.篇名.报纸名称，日期，版次。

举例如下：

\noindent{\zihaosi\bfseries 参考文献（四号、黑体、顶格）}

\begingroup
\zihaowu
\parindent 0pt
\newcommand{\NKTmanualref}[1]{\hangindent=2.4em\hangafter=1\noindent #1\par}
\NKTmanualref{[1] 庞青山.论大学学科组织及其特色.高等理科教育，2005，63（5）：1\textasciitilde3.}
\NKTmanualref{[2] Koh Y W, Lai C S, Loh K, \textit{et al}. Growth of bismuth sulfide mamowire using bismuth trisxanthate single sourceprecursors. Chem Mater, 2003, 15(24): 4544\textasciitilde4554.}
\NKTmanualref{[3] 李明.物理学.北京：科学出版社，1977，58\textasciitilde62.}
\NKTmanualref{[4] Dupont B. Bone marrow transplantation in severe combined immunodeficiency with an unrelated MLC compatible donor. In：White H J, Smith R, eds. Proceedings of the Third Annual Meeting of the International Society for Experimental Hematology. Houston：International Society for Experimental Hematology, 1974.44\textasciitilde46．}
\NKTmanualref{[5] 胡刚.蛋白质深度分析以及基因的进化模型：[博士学位论文].天津:南开大学，2005.}
\NKTmanualref{[6] 姚光起.一种氧气镐材料的制备方法.中国专利.ZL891056088，1980-07-03.}
\NKTmanualref{[7] 中华人民共和国国家技术监督局.GB3100-3102.中华人民共和国国家标准.北京：中国标准出版社，1994-11-01.}
\endgroup

\textcolor{red}{以上序号用中括号，与文字之间空两格。如果需要两行的，第二行文字要位于序号的后边，与第一行文字对齐。中文用五号宋体，外文用五号 Times New Roman 字体。}

\section{致谢}

致谢是作者对该文章的形成作过贡献的组织或个人予以感谢的文字记载，语言要诚恳、恰当、简短。

\section{附录}

是否需要附录可根据毕业论文（设计）情况而定。附录应另起一页，内容一般包括正文中不便列出的冗长公式推导、符号说明（含缩写）、计算机程序等。\textcolor{red}{“附”“录”中间空两格、四号字、黑体、居中，内容采用小四号、宋体。}

\chapter{书写要求} \label{chpt:writing}

\section{文字、标点符号和数字}

学位论文一律用汉字书写。除特殊需要外，不得使用已废除的繁体字、异体字等不规范汉字。标点符号的用法应该以 GB/T 15834—1995《标点符号用法》为准。

\section{层次标题}

毕业论文（设计）正文数字标题书写顺序依次为：一、（一）1. (1) ①；第一级标题用小三号黑体字，第二级标题用四号黑体字，第三级及以下标题用小四号黑体字。

\section{篇眉和页码}

\textcolor{red}{页眉从中文摘要开始，按章进行更换，内容与本章的标题相同；论文页眉奇偶页相同，论文各章的首页也有篇眉。}论文首页以及第二页（中文摘要之前的部分）不需要页眉。

页码从正文的第 1 章（引言）开始按阿拉伯数字（1，2，3……）连续编排，之前的部分（中文摘要、Abstract、目录等）用大写罗马数字（Ⅰ，Ⅱ，Ⅲ……）单独编排。

\textcolor{red}{第一章  ××××（行居中）}

\textcolor{red}{第一节  ××××（行居中）}

\textcolor{red}{1.1.1  ××××（行左首）}

\section{有关图、表、表达式}

\subsection{图}

a.毕业论文（设计）的插图必须精心制作，线条要匀洁美观。插图应与正文呼应，不得与正文无关或与正文脱节；正文中要求对插图进行解释说明。

b.图的内容安排要适当，不要过于密实。

c.每幅插图应有题目和序号，全文的插图可以统一编序，也可以逐章单独编序，如：图 45 或图 6.8；采取哪一种方式应和表格、公式的编序方式统一。图序必须连续，不重复，不跳缺。

d.由若干分图组成的插图，分图用 a、b、c……标序。分图的图名以及图中各种代号的意义，以图注形式写在图题下方，先写分图名，另起行后写代号的意义。例：

\vspace{7\baselineskip}
\begin{center}
图 2.2\quad 纹理块的边界
\end{center}

\subsection{表格}

a.表格必须与论文叙述有直接联系，不得出现与论文叙述脱节的表格。表格中的内容在技术上不得与正文矛盾。

\textcolor{red}{b.每个表格都应有表题和序号。表题应写在表格上方正中，序号写在左方，不加标点，空一格接写标题，表题末尾不加标点。}

\textcolor{red}{c.全文的表格可以统一编序，也可以逐章单独编序。采用哪一种方式和插图、公式的编序方式统一。表序必须连续，不得跳缺。正文中引用时，“表”字在前，序号在后，如写“表 2”。}

d.表格允许下页接写，接写时表题省略，表头应重复书写，并在右上方写“续表××”。多项大表可以分割成块，多页书写，接口处必须注明“接下页”、“接上页”、“接第×页”字样。

\textcolor{red}{e.表格采用“三线表”形式，应位于正文首次出现处的段落下方，不应置前和过分置后。}

例：

\begin{center}
\zihaowu
表 3.1\quad 直接和间接方式的比较

\vspace{0.5\baselineskip}
\begin{tabular}{lcc}
\toprule
 & 间接方式 & 直接方式 \\
\midrule
网络效率 & 低 & 高 \\
呼叫建立迟延 & 无 & 无 \\
呼叫处理迟延 & 轻 & 轻 \\
网络复杂性 & 简单 & 复杂 \\
\bottomrule
\end{tabular}
\end{center}

\subsection{表达式}

a.公式应另起一行写在稿纸中央。一行写不完的长公式，最好在等号处转行，如做不到这一点，可在数学符号（如“+”、“-”号）处转行。

b.公式的编号用圆括号括起，放在公式右边行末，在公式和编号之间不加虚线。公式可按全文统编序号，也可按章独立序号，如（49）或（4.11），采用哪一种序号应和图序、表序编法一致。不应出现某章里的公式编序号，有的则不编序号。子公式可不编序号，需要引用时可加编 a、b、c……重复引用的公式不得另编新序号。公式序号必须连续，不得重复或跳缺。

c.文中引用某一公式时，写成“由式（16.20）”。

\section{参考文献}

\textcolor{red}{每篇文献最多列出 3 位作者，超出 3 位时，中文写“等”，英文写“et al”（斜体）。作者的姓名一律姓在前名在后，欧美人的名字可以用缩写字母，且缩写名后省略缩写点“.”。参考文献应另起一页，一律放在正文后。参考文献应按文中引用出现的顺序排列。正文中引用参考文献的部位，须用上标标注[文献号]。}

\chapter{排版及印刷要求} \label{chpt:typeset}

\section{纸张要求及页面设置}

\textcolor{red}{有关纸张要求及页面设置等有关要求，请参见表 4.1。}

{\color{red}\begin{center}
\zihaowu 表4.1\quad 纸张要求及页面设置
\end{center}}

\begin{center}
\renewcommand{\arraystretch}{1.25}
\zihaowu
\begin{tabularx}{0.86\textwidth}{|p{3cm}|X|}
\hline
要\qquad 求 & \\
\hline
纸张 & A4（210×297），幅面白色 \\
\hline
页面设置 & 上、下 2.54cm，左、右 3.17cm，页眉 1.5cm，页脚 1.75cm（默认设置） \\
\hline
\end{tabularx}
\end{center}

\textcolor{red}{注意：本科生毕业论文要求使用“三线表”(见第 7 页)，切记}

\section{中文封面（已经定型，只需填入相关内容）}

\textcolor{red}{直接根据范文上面的首页进行修改即可。}

\section{中、英文摘要}

{\color{red}\begin{center}
\zihaowu 表 4.3\quad 中、英文摘要格式要求
\end{center}}

\begin{center}
\renewcommand{\arraystretch}{1.25}
\zihaowu
\begin{tabularx}{0.92\textwidth}{|p{2.2cm}|X|X|}
\hline
 & 中文摘要 & 英文摘要 \\
\hline
标题 & 摘要：“摘”“要”中间空两格、四号字、黑体、居中 & Abstract：四号 Times New Roman、加粗、居中 \\
\hline
关键词 & “关键词”三个字用小四号字、黑体、顶格写 & “Keywords”用小四号字、Times New Roman、加粗、顶格写 \\
\hline
\end{tabularx}
\end{center}

\section{目录}

{\color{red}\begin{center}
\zihaowu 表 4.4\quad 目录格式要求
\end{center}}

\begin{center}
\renewcommand{\arraystretch}{1.25}
\zihaowu
\begin{tabularx}{0.9\textwidth}{|p{4.6cm}|X|}
\hline
示例 & 要\qquad 求 \\
\hline
标题\quad 目录 & “目”“录”中间空两格、三号字、黑体、居中 \\
\hline
各章目录\quad 第一章\quad 格式要求„„„„„„5 & 每章题目采用小三号宋体字 \\
\hline
节标题目录\quad 第一节\quad 有关图、表、表达式„„10 & 每节题目采用四号宋体字 \\
\hline
\end{tabularx}
\end{center}

\newpage

{\color{red}\begin{center}
\zihaowu 续表4.4\quad 目录格式要求
\end{center}}

\begin{center}
\renewcommand{\arraystretch}{1.25}
\zihaowu
\begin{tabularx}{0.9\textwidth}{|p{4.6cm}|X|}
\hline
小节标题目录\quad 3.5.1\quad 图„„„„„„„„„10 & 每节题目采用小四号宋体字 \\
\hline
\end{tabularx}
\end{center}

\textcolor{red}{注意：至多生成三级目录，两级目录也是可以的。}

\section{正文}

\textcolor{red}{正文部分是论文的核心，格式要求如表 4.5 所示，请大家认真研读，特别需要注意一下问题：1）孤行问题：章节标题不能在一页的最后一行，下面需要有内容。2）图表要有编号和名称，编号和名称要和图表在同一页。3）由于图表的原因，可能产生较大篇幅的空白页，需要调整图表的位置。4）当表较大时，需要跨页，每页需要重复表头，加上表名-“续表 X.X\ \ XXXXX”。5）表的内容格式要求：宋体，10 磅，单倍行距，段前 0 磅，段后 0 磅。}

{\color{red}\begin{center}
\zihaowu 表 4.5\quad 正文格式要求
\end{center}}

\begin{center}
\renewcommand{\arraystretch}{1.25}
\zihaowu
\begin{tabularx}{0.92\textwidth}{|p{4cm}|X|}
\hline
示例 & 要求 \\
\hline
一级标题\quad 第一章\quad ×× & 黑体小三（行居中） \\
\hline
二级标题\quad 第一节\quad ××× & 黑体四号（行居中） \\
\hline
三级及以下\quad 1.1.1\quad ××× & 黑体小四（行左端） \\
\hline
段落文字\par ××××××××××\par ××××××××××××××\par ×××××××××××××× & 宋体小四（英文用 Times New Roman 体 12 磅），两端对齐书写，段落首行左缩进 2 个汉字符，1.5 倍行距。 \\
\hline
图序、图名\quad 图 2.1\quad ××× & 置于图的下方，宋体 5 号 \\
\hline
表序、表名\quad 表 3.1\quad ××× & 标题写在表格上方正中，序号空一格写标题，宋体 5 号 \\
\hline
表达式\quad （3.2） & 序号加圆括号，放在公式右边行末 \\
\hline
\end{tabularx}
\end{center}

\section{程序部分}

\textcolor{red}{一般毕业论文中不允许出现大量代码。如果必需出现，英文采用 Times New Roman 字体，汉字用宋体，所有代码的字号应该比正文字体要小，用“五号”字。程序行间距应该为单倍间距。程序部分底色应该为白色。}

\section{其他}

\begin{center}
\renewcommand{\arraystretch}{1.25}
\zihaowu
\begin{tabularx}{0.92\textwidth}{|p{2.4cm}|X|}
\hline
项目 & 要求 \\
\hline
注释 & 一律采用页下注（脚注）注释内容当页完成，中文用小五号宋体，英文用小五号 Times New Roman 字体。注释序号用①②③，与注释文字之间空一格。在同一页中有两个及以上的注释时，按注释在正文中的先后顺序编号，并标注在正文右上角，如×××①。每一页独立编号。 \\
\hline
参考文献 & 标题要求同第一级标题，内容格式如文档给出 \\
\hline
附录 & “附”“录”中间空两格、四号字、黑体、居中，内容采用小四号、宋体。 \\
\hline
致谢 & “致谢”二字中间空两格、四号字、黑体、居中。内容限 1 页，采用小四号宋体。 \\
\hline
\end{tabularx}
\end{center}

\section{印刷及装订要求}

毕业论文（设计）结构一般由封面、中英文内容摘要及关键词、目录、正文、附录、参考文献、致谢等组成。毕业论文（设计）一律采用打印形式，编排装订顺序依次为：（1）封面（2）声明（3）中英文内容摘要及关键词（4）目录（5）正文（6）附录（7）参考文献（8）致谢（9）毕业论文（设计）题目审批表（10）毕业论文（设计）中期检查表（11）毕业论文（设计）指导教师评语及打分表（12）毕业论文（设计）答辩记录及打分表（13）“查重”检测结果及对检测结果的认定材料。

\textcolor{red}{字符间距为默认值（缩放 100\%，间距：标准）；论文总页数在 50 页以上的要求双面打印。}

\chapter{总结与展望} \label{chpt:conclusion}

\section{主要工作}

\textcolor{red}{论文题目应简明扼要地概括和反映出论文的核心内容，要求在 20 字以内；}论文题目应具有高度的概括性，且简明、易读。

学位论文形式结构主体部分包括：引言，正文，结束语；论文一定要重点突出，总体上要求：图表清晰，叙述流畅，版面整洁，章节有序，层次分明。

论文陈述的是客观事实，不是主观的情感，在用词方面不能用“我”、“本人”、“作者”、“笔者”，要用第三人称，至多用“我们”，可以用“本文如何如何”、“系统如何如何”。

最后，强调论文不是系统说明书或技术文档；不可出现过多代码，大段代码可放在附录，不要过多介绍现有知识、技术；章节分布要均匀；一定把自己的实现工作表述完整、表述清楚。

\section{展望}

\textcolor{red}{需要总结不足，需要改进的方面，拟进一步研究的内容和工作。}

\textcolor{red}{文字复制比（包括引用文献和本人文献部分）达到 20\%及以上的，“查重”结果首先交由指导教师复核认定，有异议的需填写《南开大学本科毕业论文（设计）“查重”结果认定表》并提交相关证据材料。“查重”结果和指导教师复核认定材料交学院学术委员会进行最终认定。}

\textcolor{red}{第二次“查重”的文字复制比低于 20\%的，可参加毕业论文（设计）答辩；第二次“查重”的文字复制比仍达到 20\%及以上的，则取消答辩资格，毕业论文（设计）成绩以 0 分计。}
